\section{Uvod}

Svet se vedno spreminja. Staro se spremeni v novo. Linearna logika poskuša modelirati to spremninjanje in porabljanje virov za doseganje cilja. 
Če imamo dva evra, lahko kupimo eno čokolado ali pa dva paketa bombonov -- nič več, nič manj. 

Linearno logiko lahko tudi uporabimo v programiranju. Oglejmo si tri funkcije na seznamih:
\begin{align*}
    & \cd{make} : \cd{Int} \rightarrow A \rightarrow A \cd{ Array} \\
    & \cd{lookup} : \cd{Int} \rightarrow A \cd{ Array} \rightarrow A \\
    & \cd{update} : \cd{Int} \rightarrow A \rightarrow A \cd{ Array} \rightarrow A \cd{ Array}
\end{align*}
Na problem naletimo pri \(\cd{update}\), ki med drugim sprejme seznam in vrne spremenjen seznam. Kaj se zgodi s starim seznamom. Še vedno obstaja, ali se je spremenil? Seznami se načeloma spreminjajo na mestu, torej naš stari seznam no obstaja več. S pomočjo linearnih tipov pa lahko zagotovimo, da do objektov ki ne obstajajo več, ne moremo dostpati.

V okviru te diplomske naloge, sem za projekt The Programming Languages Zoo~\cite{plzoo} sestavila programski jezik Line Mix, ki predstavlja minimalno različico jezika z linearnimi tipi. V tem diplomskem delu bom predstavila linearne tipe, ki so implementiratni v Line Mix, katere probleme rešujejo, in kako rešiti nove probleme, ki se ustvarijo. 


\section{Pravila sklepanja}

Teoretični del programski jezikov bom predstavila s pomočjo pravil sklepanja. Splošna oblika pravila sklepanja je~\cite{Jazbec2022Primerjava}
%
\begin{prooftree}
    \AXC{\(\mathcal{H}_1\)}
    \AXC{\(\mathcal{H}_2\)}
    \AXC{\(\ldots\)}
    \AXC{\(\mathcal{H}_n\)}
    % \RL{ime,}
    \QuaternaryInfC{\(\mathcal{S}\)}
\end{prooftree}
%
pri čemer so \(\mathcal{S}, \mathcal{H}_1, \ldots, \mathcal{H}_n\) sodbe oblike \(\G \vdash a: A\), kjer je \(\G\) kontekst, \(a\) pa je spremenljivka tipa \(A\), ki jo lahko dobimo iz tega konteksta. Poglejmo si to na primeru pravila vpeljave in uporabe za funckijski tip.


\begin{primer}
    Pravilo uporabe za funkcijo je preprosto: Če imamo funkcijo \(f : A \rightarrow B\) in element \(a : A\), potem lahko apliciramo \(a\) na \(f\) in dobimo element tipa \(B\).
    \begin{prooftree}
        \AXC{\(\G \vdash f : A \rightarrow B\)}
        \AXC{\(\G \vdash a : A\)}
        \RL{\(\rightarrow\)E}
        \BIC{\(\G \vdash f~a : B\)}
    \end{prooftree}
    Pravilo vpeljave za funkcijo pa se glasu: Če lahko iz konteksta \(\G\) z elementom \(x: A\) dobimo \(b : B\), potem lahko iz konteksta \(\G\) dobimo funkcijo \(\la x : A. B : A \rightarrow B\). 

    \begin{prooftree}
        \AXC{\(\G, x: A \vdash b: B\)}
        \RL{\(\rightarrow\)I}
        \UIC{\(\G \vdash \la x : A. b : A \rightarrow B\)}
    \end{prooftree}
\end{primer}
